\partdivider{E}{Gravitation}{The volume's finale turns every tool skyward. One inverse-square force, and the conservation laws built earlier, explain falling apples, orbiting moons, and all three of Kepler's laws.}{\chref{ch:gravitation} Gravitation: Force, Field and Energy\\[1.2mm]
\chref{ch:kepler} Kepler's Laws and Orbits}

\ncertchapter{16}{Gravitation: Force, Field and Energy}\label{ch:gravitation}

Every force met so far in this volume acts by contact: a push, a pull through a string, the grip of
friction. Gravity is different. It reaches across empty space, between bodies that never touch, and it
acts on everything with mass in exactly the same way. This chapter applies the whole machinery of the
volume (vectors, Newton's laws, energy, and integration over a body sliced into pieces) to that single
force, and the result is an account of why apples fall, why the Moon does not, and why a rocket needs a
particular speed to leave the Earth for good.

The force law is stated at once; its consequences repay slow derivation. The shell theorem explains why a
planet attracts as though all its mass were at its centre, and why, in a uniform planet, gravity is
strongest at the surface. The potential energy, measured from a zero at infinity, turns out to be
negative, and the sign carries physics: the Earth and a body near it are bound together. The escape speed
follows from setting the total energy to zero.

The chapter ends by being candid about the law's limits. Newton's gravity fails, slightly but measurably,
in the orbit of Mercury, the bending of starlight and the running of clocks, and the final section
explains how Einstein's general relativity replaces the force with the curvature of spacetime, reducing
to Newton's law in the weak, slow limit.

\section{Newton's Law of Universal Gravitation}\label{sec:newtonlaw}

The first question is the force itself: how strongly do two masses attract, and how does the pull
depend on the distance between them? Newton's answer is a single statement about any two particles in
the universe.

\begin{keybox}{Newton's law of gravitation}
Every particle attracts every other with a force along the line joining them. For particles of masses
$m_1$ and $m_2$ a distance $r$ apart, with $\uhat{r}$ the unit vector from $m_1$ towards $m_2$, the force
on $m_2$ is
\[
\vF=-\frac{Gm_1m_2}{r^2}\,\uhat{r},\qquad F=\frac{Gm_1m_2}{r^2},\qquad G=\SI{6.674e-11}{N.m^2/kg^2}.
\]
Here $\vF$ is a vector (in newtons) and $F$ its magnitude; $G$ is the \emph{universal gravitational
constant}. The force is always attractive, obeys the third law exactly, and is a \emph{central} force
(it acts along the line joining the two bodies), the property on which the whole of \chref{ch:kepler} rests.
\end{keybox}

The word doing the work in the title is \emph{universal}. Before Newton, the heavens and the Earth were
different places as a matter of physics: terrestrial things fell, celestial things circled eternally,
and few doubted that they obeyed different laws. Newton's claim was that the force pulling an apple down is
\emph{the same force}, with the same constant $G$, holding the Moon in orbit, differing only in $r$
(Figure~\ref{fig:gravuniv} shows the check Newton made).

Three remarks make the law usable. First, it is stated for \emph{particles}, bodies small compared with
their separation; for extended bodies one must add up the pulls of all their pieces, and the next section
shows that for spheres this sum comes out exactly as if each were a point at its centre. Second, that
adding up is legitimate because gravity obeys the \emph{principle of superposition}: the force between
two particles is unaffected by the presence of others, so the net force on a body is the \emph{vector}
sum of the separate pulls,
\[
\vF_{\text{net}}=\vF_1+\vF_2+\vF_3+\cdots .
\]
Third, $G$ is tiny. Two \SI{50}{kg} students a metre apart attract each other with a force
of only $Gm_1m_2/r^2\approx\SI{1.7e-7}{N}$, about the weight of a fine grain of
sand, which is
why gravity between everyday objects goes unnoticed and only a body of planetary mass makes it obvious.
By the third law an apple pulls the Earth up exactly as hard as the Earth pulls the apple
down; the Earth does not visibly move because the same force divided by \SI{6e24}{kg} gives an
acceleration of order $10^{-25}$\,\si{m/s^2}.

\begin{figure}[htbp]
\centering
\tikzfig{grav_universal}
\figcaption{fig:gravuniv}{The evidence for ``universal''. The Moon sits $60$ Earth-radii from the centre,
so an inverse-square law predicts it feels $1/60^2=1/3600$ of surface gravity. Measuring the Moon's
period and computing its centripetal acceleration gives $\SI{0.0027}{m/s^2}$, which is $g/3600$ to
within the accuracy of these rounded figures: one number from a stopwatch pointed at the sky, the other
from dropping things, and nothing connects them unless the same law runs both.}
\end{figure}

\begin{trapbox}{$g$ and $G$ are not the same kind of thing}
$G$ is a \emph{constant of nature}: the same everywhere, dimensions $\mathsf M^{-1}\mathsf L^3\mathsf T^{-2}$,
depending on nothing. $g$, the acceleration of free fall met since \chref{ch:straightline}, is a
\emph{local fact} about the body you stand on. Setting a body's weight $mg$ equal to $GMm/R^2$, with $M$
and $R$ the planet's mass and radius (\S\ref{sec:shell} justifies treating the planet as a point at its
centre), gives $g=GM/R^2$: \SI{9.8}{m/s^2} on
Earth, \SI{1.6}{m/s^2} on the Moon, dimensions $\mathsf L\mathsf T^{-2}$. Confusing them is this chapter's
commonest error. Note that measuring $g$ gives only the \emph{product} $GM$, never $G$ alone, which is
why weighing the Earth needs an independent measurement of $G$ (Cavendish's torsion balance, 1798).
\end{trapbox}

\subsection*{A note on the discovery of Neptune}
\addcontentsline{toc}{subsection}{A note on the discovery of Neptune}

\begin{remarkbox}{A planet found with a pen before it was found with a telescope}
The strongest single piece of evidence that Newton's law is \emph{universal} is not an experiment at all.
It is a planet that was located by arithmetic.

\smallskip
\textbf{The anomaly.} Uranus was discovered by William Herschel in 1781. Once its orbit had been fitted,
astronomers could predict where it should be on any future night. By the 1820s the predictions were
failing. The discrepancy was tiny, around two arcminutes, about one-fifteenth of the Moon's apparent
diameter, but it was far larger than the observational error and it would not go away. Either Newton's
law was wrong at that distance, or something not yet seen was pulling on Uranus.

\smallskip
\textbf{The inversion.} Up to that point gravitation had always been used \emph{forwards}: given the
masses and positions, compute the motion. Two astronomers ran it \emph{backwards}: given the error in
the motion, compute the mass and position of the body responsible. John Couch Adams in Cambridge and
Urbain Le Verrier in Paris did this independently, each assuming an unseen eighth planet and solving for
the orbit that would produce exactly the observed deviation. This is an inverse problem, and it is a far
harder calculation than the forward one, because many combinations of mass and distance give nearly the
same pull.

\smallskip
\textbf{The test.} Le Verrier sent his predicted position to Johann Galle at the Berlin Observatory. On
23~September~1846, the night the letter arrived, Galle and his student Heinrich d'Arrest compared the sky
at those coordinates with a new star chart and within an hour found an object that was not on it, within
$1^{\circ}$ of the prediction. (Adams' less precise predictions led James Challis at Cambridge to record
the planet twice that August without recognising it.)

\smallskip
\textbf{Why it matters for this chapter.} A theory that merely \emph{accounts for} what we already see can
always be patched. A theory that says ``point a telescope at this coordinate on this night and you will
find a new world'' is risking everything on one number, and $F=Gm_1m_2/r^{2}$ survived that risk. Notice
also which pieces of this chapter the calculation used: the inverse-square law, the superposition of the
pulls of several bodies, and the treatment of each planet as a point mass at its centre, which is the
shell theorem of \S\ref{sec:shell}.
\end{remarkbox}

\why{\textbf{The same method, applied again, failed, and that failure built the next theory.}
Mercury's orbit showed an anomaly too: the point of closest approach to the Sun creeps forward by about
\ang{;43} per century more than the other planets can account for. Le~Verrier, fresh from Neptune,
applied his method a second time in 1859 and proposed an unseen planet inside Mercury's orbit, soon named
Vulcan. Vulcan does not exist, and decades of searching never found it. The residue was not another mass; it was a flaw
in the law itself, and explaining those \ang{;43} per century became one of the first triumphs of
Einstein's theory of gravity. That story is taken up in \S\ref{sec:beyondnewton}.}

\section{\texorpdfstring{The Shell Theorem, and $g$ Away from the Surface}{The Shell Theorem, and g Away from the Surface}}\label{sec:shell}

The law speaks of particles, yet the Earth is anything but small compared with the distance to an apple
on its surface, so before the law can be used for planets we must know how a large sphere pulls. Every
formula in this chapter treats a sphere as a point at its centre. For a spherically symmetric
body that is not an approximation but an exact result, and it rests on the shell theorem. It is stated
here and proved in \S\ref{sec:shellproof}, once the potential has made the proof short.

\begin{keybox}{The shell theorem}
A uniform spherical shell attracts an external body \emph{as if all its mass were at its centre}. And
\emph{inside} a uniform shell, the net gravitational force on a body is exactly zero, everywhere. (Both results depend
on the force being inverse-square, which is why Coulomb's law in electrostatics will feel like a rerun.)
A solid sphere whose density depends only on the distance from the centre is a nest of such shells, so
it too attracts every outside body as if all its mass were at its centre.
\end{keybox}

The inside result is the more surprising one, and there is a clean intuition for it. Stand off-centre
inside a hollow shell and draw a narrow double cone through your position, opening equally in opposite
directions. The patch it cuts from the near wall is close, so each kilogram of it pulls hard, but the
patch is small. The patch it cuts from the far wall pulls more weakly per kilogram but is larger. The
area, and hence the mass, of each patch grows as the square of its distance from you ($\propto r^2$),
while the pull per unit mass falls as the inverse square ($\propto 1/r^2$). (Both patches are also tilted
by the same angle to the cone, because a chord meets a sphere at equal angles at its two ends.) The two
$r^2$ factors cancel \emph{exactly}, so the near and far pulls balance, and they do so for every cone
at once. This precise cancellation
happens \emph{only} for an inverse-square law; any other power would leave a residual field inside.

The theorem's first payoff is the answer to a practical question: how does $g$ change as you climb above
the surface, or go down a mine below it?

\begin{derivbox}{How $g$ varies with altitude and depth}
Model the Earth as a uniform sphere of mass $M$, radius $R$, with surface gravity $g=GM/R^2$.
\stepn{1} \textbf{Above the surface}, at height $h$, the whole Earth still acts as a point at its centre,
now a distance $R+h$ away:
\[
g_h=\frac{GM}{(R+h)^2}=g\left(1+\frac{h}{R}\right)^{-2}\approx g\left(1-\frac{2h}{R}\right)\ \text{for }h\ll R.
\]
The binomial approximation is only valid for small $h$; for a satellite it is badly wrong, and the exact
$(R+h)^{-2}$ must be used.
\stepn{2} \textbf{Below the surface}, at depth $d$, the shell theorem does real work: the shell of
material \emph{above} you (farther from the centre) contributes exactly \emph{zero}, and only the inner
sphere of radius $R-d$ pulls. For a uniform Earth its mass scales as the cube of the radius,
$M'=M\left(\frac{R-d}{R}\right)^3$, and it acts from distance $R-d$:
\[
g_d=\frac{GM'}{(R-d)^2}=\frac{GM}{R^2}\cdot\frac{R-d}{R}=g\left(1-\frac{d}{R}\right).
\]
So $g$ falls \emph{linearly} to zero at the centre, where you would be pulled equally in all
directions, hence weightless.
\stepn{3} \textbf{The asymmetry.} Rising costs you $2h/R$ (to first order), but descending only $d/R$:
gravity falls off \emph{twice as fast going up as going down}, and for a uniform Earth the surface is
exactly where $g$ is largest (Figure~\ref{fig:gravshell}).
\end{derivbox}

\begin{remarkbox}{Where the uniform-Earth model fails}
The real Earth is not uniform: its iron core is more than twice as dense as the rocks above it, so
$M'$ does not fall as $(R-d)^3$ near the surface. Going down through the low-density crust and mantle
removes little mass while bringing you closer to the dense core, and the measured $g$ actually
\emph{rises} slightly, to about \SI{10.7}{m/s^2} at the top of the core, \SI{2900}{km} down, before
falling to zero at the centre. The shell theorem is untouched (only the mass inside $r$ still counts);
it is the assumption of uniform density that fails. Unless told otherwise, examination problems assume
the uniform model.
\end{remarkbox}

\begin{figure}[htbp]
\centering
\tikzfig{grav_shell}
\figcaption{fig:gravshell}{For a uniform planet, gravity rises linearly from zero at the centre to its
maximum at the surface, then falls off as $1/r^2$ outside. The surface is the peak, and the corner there
is real: the two laws meet but have different slopes. Inside, only the mass closer to the centre than
you counts; the shell above contributes nothing.}
\end{figure}

\begin{exbox}{Three questions, one equation}
\textbf{(i) At what altitude is $g$ half its surface value?} Set $\dfrac{GM}{(R+h)^2}=\dfrac12\dfrac{GM}{R^2}$,
so $(R+h)^2=2R^2$, giving $h=R(\sqrt2-1)=0.414R=\SI{2650}{km}$. You must climb 41\% of the Earth's radius
to halve your weight; Everest, at \SI{8.8}{km}, changes $g$ by only 0.3\%. (The small-$h$
formula would give $0.25R$, off by 40\%, because $h\ll R$ is violated.)
\textbf{(ii) Mars} ($M=0.107M_E$, $R=0.532R_E$): using $g=GM/R^2$ as a \emph{ratio} so $G$ cancels,
$\dfrac{g_M}{g_E}=\dfrac{0.107}{0.532^2}=0.378$, so $g_M=\SI{3.7}{m/s^2}$, about $\tfrac{3}{8}$ of Earth's.
\textbf{(iii) The ISS} at \SI{400}{km}: $g_h=g\left(\dfrac{6400}{6800}\right)^2=\SI{8.7}{m/s^2}$, which is
89\% of surface gravity. And yet the astronauts float, because they are in free fall, not because
gravity is absent (the freely falling lift of \chref{ch:newton}). ``Zero gravity'' is a misnomer.
\ans{\SI{2650}{km}; \SI{3.7}{m/s^2}; \SI{8.7}{m/s^2}}
\end{exbox}

\section{The Gravitational Field}\label{sec:field}

The results so far have quietly changed point of view. A formula such as $g_h=GM/(R+h)^2$ says
nothing about the body that falls; it describes what the Earth does at each point of the space around
it, whatever is placed there. It is often more useful to work with that description than to compute
the force on a particular body, and it has a name: the gravitational field.

\begin{defbox}{Gravitational field}
The gravitational field $\vc{E}_g$ at a point is the force per unit mass that a small test mass $m$
placed there would feel, a \emph{vector} property of the point in space, defined whether or not a test
mass is present to feel it:
\[
\vc{E}_g=\frac{\vF}{m}=-\frac{GM}{r^2}\,\uhat{r},\qquad\text{units }\si{N/kg}=\si{m/s^2},
\]
for the field of a point mass (or sphere) $M$ at distance $r$ from its centre, with $\uhat{r}$ pointing
away from $M$: the minus sign says the field points \emph{towards} the mass. A body of any mass placed
in the field feels $\vF=m\vc{E}_g$ and, if nothing else acts, accelerates at $\vc{E}_g$. At the
Earth's surface, $\vc{E}_g$ is just $\vg$: the ``acceleration due to gravity'' and the field strength
are the same quantity.
\end{defbox}

Picture the field as an arrow attached to every point of space, pointing at the Earth's centre and
shrinking as $1/r^2$ with distance; lines drawn along these arrows (field lines) all run radially inward
and end on the Earth. Because forces superpose, so do fields: the field of several masses is the vector
sum of their separate fields.

\begin{exbox}{Where the Earth and the Moon cancel}
The Earth's mass $M_E$ is $81$ times the Moon's, and their centres are $d=\SI{3.84e5}{km}$ apart. Find the point
on the line joining them where the net gravitational field is zero.

\stepn{1} \textbf{Superpose.} Between the two bodies the fields point in opposite directions, so they
cancel where their magnitudes are equal. At distance $x$ from the Earth's centre,
\[
\frac{GM_E}{x^2}=\frac{GM_E/81}{(d-x)^2}\quad\Longrightarrow\quad\frac{x}{d-x}=9 .
\]
\stepn{2} \textbf{Solve.} $x=0.9\,d\approx\SI{3.46e5}{km}$ from the Earth, nine-tenths of the way to the
Moon. (Outside the segment the two fields point the same way and cannot cancel.) \ans{$0.9\,d$ from the Earth}
\end{exbox}

\why{The field looks like mere notation (divide by $m$ and rename), but it relocates the physics.
Instead of ``the Sun pulls the Earth across empty space'', the Sun fills space with a field, and the Earth
responds to the field where the Earth is. In Newtonian gravity the two descriptions predict the same
things, but the field turns out to be real: changes in it travel outward at the speed of light, as the
gravitational waves of \S\ref{sec:beyondnewton}.}

\section{Gravitational Potential Energy}\label{sec:gpe}

Forces and fields answer questions about acceleration. Many gravitational questions are instead about
speeds at two different places (how fast must a rocket leave the ground to reach a given height?), and
\chref{ch:energy} showed that such questions are answered most easily by energy conservation, which
needs a potential energy. Near the Earth's surface the potential energy of a body is $mgh$. That
formula assumes a constant $g$, which fails over distances comparable to the Earth's radius. The general
expression follows from integrating the inverse-square force.

\begin{derivbox}{Why $U=-GMm/r$}
Potential energy is defined as the negative of the work done by the (conservative) gravitational force:
$U(r)=-W_{\text{grav}}$, with the reference $U=0$ at infinity.
\stepn{1} \textbf{Work done by gravity bringing $m$ inward from infinity to $r$.} At distance $x$ the
gravitational force has magnitude $GMm/x^2$ and points \emph{inward} (towards $M$), the same direction the
mass moves as it comes in, so gravity does \emph{positive} work. Integrating the inward force over the
inward displacement:
\[
W_{\text{grav}}=\int_{\infty}^{r}\frac{GMm}{x^2}\,(-dx)=GMm\left[\frac{1}{x}\right]_{\infty}^{r}=\frac{GMm}{r}>0.
\]
(The inward displacement is $-dx$, which is positive because $x$ \emph{decreases} as the mass moves in;
the two minus signs combine to give the positive result expected for an attractive force pulling
something in.)
\stepn{2} \textbf{Hence the potential energy:}
\[
U(r)=-W_{\text{grav}}=-\frac{GMm}{r}.
\]
It is \emph{negative everywhere}: we set $U=0$ at infinity, and since gravity does positive work drawing
the mass in, $U$ must fall below zero.
\end{derivbox}

\begin{remarkbox}{Read the minus sign as ``bound''}
Only \emph{differences} of $U$ are physical, and placing the zero at infinity is a choice, made because
the force vanishes there. Once that choice is made, the sign means something: $U=-GMm/r<0$ says the pair
is \emph{trapped}, so you must supply energy to separate them, and the more negative $U$, the more
tightly bound. Note also that $U$ belongs to the \emph{pair}, Earth and body together, not to the body
alone. For three or more particles the potential energy is the sum of $-Gm_im_j/r_{ij}$ over every pair;
three equal masses $m$ at the corners of an equilateral triangle of side $a$ have $U=-3Gm^2/a$.
\end{remarkbox}

\begin{derivbox}{Where $mgh$ comes from, and where it fails}
Raise $m$ slowly from the surface ($r=R$) to height $h$; the work needed equals the rise in potential
energy, which, using $GM=gR^2$ (from $g=GM/R^2$), is
\[
\Delta U=-\frac{GMm}{R+h}+\frac{GMm}{R}=GMm\,\frac{h}{R(R+h)}=\frac{mgh}{1+h/R}.
\]
For $h\ll R$ the denominator is $1$ and $\Delta U=mgh$, the result of \chref{ch:energy}: $mgh$ is the
small-height limit of $-GMm/r$, not a different law. For large $h$ the exact $\Delta U$ is always
\emph{less} than $mgh$, because $g$ weakens on the way up; at $h=R$ it is $\tfrac12mgR$, half of what
$mgh$ predicts.
\end{derivbox}

Just as the field strips the test mass out of the force, it pays to strip it out of the potential energy,
leaving a property of the point alone.

\begin{defbox}{Gravitational potential}
The gravitational potential $V$ at a point is the potential energy per unit mass of a test mass placed
there; equivalently, the work an external agent must do per unit mass to bring a test mass slowly from
infinity to that point. It is a \emph{scalar}, measured in \si{J/kg}. For a point mass or sphere $M$,
\[
V=\frac{U}{m}=-\frac{GM}{r},
\]
and the potential of several masses is the plain (not vector) sum of their separate potentials. A mass
$m$ placed where the potential is $V$ has potential energy $U=mV$.
\end{defbox}

The potential also gives the field back. Since $U=-W_{\text{grav}}$, a small outward step $dr$ changes
$U$ by $dU=-F_r\,dr$, where $F_r$ is the outward component of the force, so $F_r=-dU/dr$; per unit mass,
\[
E_r=-\frac{dV}{dr}.
\]
Check: $V=-GM/r$ gives $E_r=-GM/r^2$, inward with the right size. The field therefore points
\emph{downhill}, towards lower potential, just as a ball rolls towards lower ground. This relation is what makes the
potential so useful: it is a \emph{scalar}, so the potentials of many pieces simply add, and the field
follows by one differentiation. The next subsection uses exactly this to prove the shell theorem.

\subsection{Proof of the shell theorem}\label{sec:shellproof}

The shell theorem was stated in \S\ref{sec:shell} and has been used ever since; with the potential in
hand it can now be proved. The plan is to add up potentials, which are scalars, rather than forces,
which are vectors, and then recover the field by differentiating.

\begin{figure}[htbp]
\centering
\tikzfig{grav_ring}
\figcaption{fig:gravring}{Slicing a thin shell into rings about the line $OP$. Every point of the ring
at angle $\theta$ is the same distance $s$ from $P$, so the ring's potential at $P$ needs no
components. Notice that $s$ ranges over $2R$ when $P$ is outside the shell but only over $2r$ when $P$
is inside: that difference is the whole theorem.}
\end{figure}

\begin{derivbox}{The potential of a thin uniform shell}
Take a shell of mass $M$ and radius $R$, and a point $P$ at distance $r$ from its centre $O$
(Figure~\ref{fig:gravring}).
\stepn{1} \textbf{Slice the shell into rings} about the line $OP$. The ring at angle $\theta$ from $OP$
(measured at $O$) has radius $R\sin\theta$ and width $R\,d\theta$, so its share of the shell's area
$4\pi R^2$, and hence of its mass, is
\[
dM=M\,\frac{2\pi R\sin\theta\cdot R\,d\theta}{4\pi R^2}=\frac{M}{2}\sin\theta\,d\theta .
\]
\stepn{2} \textbf{Every point of a ring is the same distance $s$ from $P$}, so the ring's potential at $P$
is simply $-G\,dM/s$. (This is why we use the potential: a vector sum over the ring would need
components; a scalar sum needs none.) By the cosine rule,
\[
s^2=R^2+r^2-2Rr\cos\theta\quad\Longrightarrow\quad 2s\,ds=2Rr\sin\theta\,d\theta ,
\]
so $\sin\theta\,d\theta=s\,ds/(Rr)$, and the ring contributes
\[
dV=-\frac{G}{s}\cdot\frac{M}{2}\cdot\frac{s\,ds}{Rr}=-\frac{GM}{2Rr}\,ds .
\]
The $s$ has cancelled, which is the inverse-square law at work.
\stepn{3} \textbf{Add up the rings} by letting $s$ run from the nearest point of the shell to the farthest.
\emph{Outside} ($r>R$): $s$ runs from $r-R$ to $r+R$, a range of $2R$, so
\[
V=-\frac{GM}{2Rr}\cdot2R=-\frac{GM}{r},
\]
exactly the potential of a point mass $M$ at $O$. By symmetry the field of the shell points along the
radius and depends only on $r$, so $E_r=-dV/dr$ gives all of it: the field outside is that of a point
mass too.
\emph{Inside} ($r<R$): $s$ runs from $R-r$ to $R+r$, a range of $2r$, so
\[
V=-\frac{GM}{2Rr}\cdot2r=-\frac{GM}{R},
\]
the same at every inside point. A potential that does not change with position has zero slope, so by
$E_r=-dV/dr$ the field inside is exactly zero. Both halves of the shell theorem are proved.
\end{derivbox}

The relation can also be run backwards, $V(r)=V(R)-\int_R^r E_r\,dr$, to find the potential inside a
uniform \emph{solid} sphere, which tells us, for instance, how fast a body dropped down a shaft would reach
the centre. There $E_r=-g\,r/R=-GMr/R^3$ (the linear law of \S\ref{sec:shell}), and
starting from $V(R)=-GM/R$ at the surface,
\[
V(r)=-\frac{GM}{R}+\frac{GM}{2R^3}\left(r^2-R^2\right)=-\frac{GM}{2R^3}\left(3R^2-r^2\right),
\]
so the potential at the centre is $-3GM/(2R)$, one and a half times its surface value. The field is zero
at the centre, but the potential is not: zero field means a flat potential, not a zero one.

\section{Escape Velocity}\label{sec:escape}

A body thrown upward fast enough never returns. Whether it returns is decided by the sign of its total
energy $E=K+U$ (a scalar, in joules; the field always carries a subscript, $\vc{E}_g$ or $E_r$, so plain
$E$ means energy from here on), which is conserved when gravity alone acts. Far from the planet $U\to0$, so any kinetic
energy the body still has there equals $E$; since $K$ can never be negative, a body with $E<0$ can never
get far away. Figure~\ref{fig:gravwell} shows this on the potential-energy curve, in the style of
\chref{ch:energy}.

\begin{keybox}{The sign of $E=K+U$ decides the fate}
\centering
\begin{tabular}{@{}cl@{}}
\toprule
$E<0$ & \textbf{bound}: an ellipse or circle; it cannot escape.\\
$E=0$ & \textbf{exactly escapes}, arriving at infinity with zero speed. A parabola.\\
$E>0$ & \textbf{escapes with speed to spare}. A hyperbola.\\
\bottomrule
\end{tabular}
\end{keybox}

\begin{figure}[htbp]
\centering
\tikzfig{grav_well}
\figcaption{fig:gravwell}{The potential well $U=-GMm/r$, with three total energies $E=K+U$. At any
$r$ the kinetic energy is the gap $K=E-U$ between the line and the curve. For $E<0$ that gap closes at
the turning point, and beyond it (dotted) $K$ would be negative, so the body is bound. For $E\ge0$ the
gap never closes and the body reaches infinity. Escape velocity is not a special rule: it is the speed
that makes $E$ exactly zero.}
\end{figure}

The shapes named in the table (ellipse, parabola, hyperbola) are proved in \S\ref{sec:ellipse}; here only
the energy matters. For a bound body the energy that must be supplied to raise $E$ to zero, and so set it
free, is $-E$, called its \emph{binding energy}; for a body at rest on the surface it is $GMm/R$. The
minimum speed for escape follows at once.

\begin{derivbox}{Escape velocity}
To \emph{just} escape means arriving at infinity with nothing left: exactly $E=0$. At the surface
(distance $R$, speed $v_e$),
\[
E=\tfrac{1}{2} mv_e^2-\frac{GMm}{R}=0\ \Rightarrow\ v_e=\sqrt{\frac{2GM}{R}}=\sqrt{2gR},
\]
using $GM=gR^2$ in the last step. For Earth, $v_e=\SI{11.2}{km/s}$. Escaping does not mean reaching a
place where gravity stops, for it never stops; it means the body keeps slowing for ever but never turns
back.

\smallskip
It is independent of the escaping body's \emph{mass} (a pebble and a battleship need the same speed,
since $m$ cancelled) and of \emph{direction} (energy is a scalar, so a sideways shot at \SI{11.2}{km/s}
escapes as well as a vertical one, provided its path does not meet the ground; air resistance is
ignored throughout). And it is \emph{not} the speed rockets travel at: they thrust continuously, which
is a different problem. Escape velocity is the one-kick speed of an unpowered projectile.
\end{derivbox}

\begin{derivbox}{Orbiting just above the surface}
For comparison, a body circling just above the surface (radius $R$, air ignored) at the
\emph{orbital speed} $v_o$ needs gravity to supply exactly the centripetal force $mv_o^2/R$
(\chref{ch:circdyn}; Figure~\ref{fig:cannon} shows both speeds as Newton imagined them):
\[
\frac{GMm}{R^2}=\frac{mv_o^2}{R}\quad\Longrightarrow\quad v_o=\sqrt{\frac{GM}{R}}=\sqrt{gR}\approx\SI{7.9}{km/s}.
\]
Comparing with $v_e=\sqrt{2gR}$ gives $v_e=\sqrt2\,v_o$: a body already in a low orbit needs only $41\%$
more speed to leave for good. In energy terms, its kinetic energy $\tfrac12mv_o^2=GMm/2R$ is exactly half
of what escape requires.
\end{derivbox}

\begin{figure}[htbp]
\centering
\tikzfig{grav_cannon}
\figcaption{fig:cannon}{Newton's cannon, fired horizontally from an exaggerated mountain, the launch point at distance $r$ from the Earth's centre. Fire harder and harder: the slow shots land, each farther away; at $\sqrt{GM/r}$ the ground curves away as fast as the shot falls and it circles for ever; faster still it traces a larger ellipse that returns to the cannon; at $\sqrt{2GM/r}$ the path opens into a parabola and it never returns. Every path drawn is an exact conic with the Earth's centre at a focus, and each slow shot lands because its closest approach to the centre lies inside the Earth.}
\end{figure}

\why{Escape velocity explains which worlds keep an atmosphere. Gas molecules have a spread of speeds, and
the fast tail leaks away if it exceeds $v_e$. The Moon ($v_e=\SI{2.4}{km/s}$) has lost almost all of its gas; Earth
($\SI{11.2}{km/s}$) kept its air but lost its hydrogen and helium, which, being lightest, move
fastest at a given temperature. This chapter and kinetic theory (a later course) jointly explain why you breathe nitrogen
and not hydrogen.}

\begin{exbox}{How high does a projectile go?}
A projectile is fired vertically from the Earth's surface at \SI{8}{km/s}. Find its greatest height
($g=\SI{9.8}{m/s^2}$, $R=\SI{6.4e6}{m}$).
(Do \emph{not} use $\tfrac{1}{2} v^2=gh$: $g$ is not constant over thousands of kilometres.) Energy conservation between the surface ($r=R$, speed $v$) and the top
($r=r_{\max}$, speed 0), using $GM=gR^2$:
\[
\tfrac{1}{2} v^2-gR=-\frac{gR^2}{r_{\max}}\ \Rightarrow\ r_{\max}=\frac{gR^2}{gR-\tfrac{1}{2} v^2}.
\]
With $gR=\SI{6.27e7}{m^2/s^2}$ and $\tfrac{1}{2} v^2=\SI{3.2e7}{m^2/s^2}$: $r_{\max}=\SI{1.31e7}{m}$, so
$h=r_{\max}-R\approx\SI{6700}{km}$, more than an Earth-radius, and about twice the \SI{3270}{km} that the
constant-$g$ formula predicts, because the projectile climbs into weaker gravity. As $\tfrac{1}{2} v^2\to gR$
the denominator vanishes and $r_{\max}\to\infty$, which \emph{is} the escape condition, reached by a
different route. \ans{$\approx\SI{6700}{km}$}
\end{exbox}

\section{Worked Examples}

These follow the order of the chapter: $g$ at, above and below the surface, then escape and orbital
energy, and finally the effect of the Earth's spin. Each states the principle before the arithmetic. Take $G=\SI{6.67e-11}{N.m^2/kg^2}$.

\begin{exbox}{Weighing the Earth}
The Earth has radius \SI{6.4e6}{m} and mass \SI{6.0e24}{kg}. Find $g$ at its surface.

\stepn{1} \textbf{Equate the weight to the gravitational force} on a test mass $m$ at the surface:
\[
mg=\frac{GM_Em}{R_E^{2}}
\qquad\Longrightarrow\qquad
g=\frac{GM_E}{R_E^{2}} .
\]
The test mass cancels, which is why $g$ is a property of the Earth and not of the object.

\stepn{2} \textbf{Substitute.}
\[
g=\frac{(6.67\times10^{-11})(6.0\times10^{24})}{(6.4\times10^{6})^{2}}
=\frac{4.0\times10^{14}}{4.1\times10^{13}}=\SI{9.8}{m/s^2}.
\]
\ans{\SI{9.8}{m/s^2}}

\smallskip
\textbf{Read it backwards.} Cavendish's torsion balance measured the attraction between lead spheres in a
laboratory, which fixes $G$; with $g$ and $R_E$ known, this equation then gives the Earth's mass, which is
why his experiment is remembered as ``weighing the Earth''.
\end{exbox}

\begin{exbox}{Height versus depth}
Find the fraction of surface gravity at (a) a height $h=R/2$, (b) a depth $d=R/2$. Then compare the
rate of fall-off near the surface.

\stepn{1} \textbf{Above the surface}, only the inverse-square law matters:
\[
g_h=\frac{gR^{2}}{(R+h)^{2}}=\frac{gR^{2}}{(1.5R)^{2}}=\frac{4g}{9}\approx0.444\,g .
\]

\stepn{2} \textbf{Below the surface}, the shell theorem says only the mass \emph{inside} the radius
pulls, giving the exactly linear law:
\[
g_d=g\left(1-\frac{d}{R}\right)=g\left(1-\tfrac{1}{2}\right)=0.5\,g .
\]

Even at half a radius the height loses more than the depth, though by less than the factor of two that
holds near the surface (\S\ref{sec:shell}).
\ans{(a) $4g/9$; (b) $g/2$}
\end{exbox}

\begin{exbox}{A planet of the same density}
A planet has twice the Earth's radius but the same mean density. How does its surface gravity compare?

\stepn{1} \textbf{Express $g$ through the density rather than the mass.} With $\rho$ the mean
density (in \si{kg/m^3}), $M=\tfrac{4}{3}\pi R^{3}\rho$, so
\[
g=\frac{GM}{R^{2}}=\frac{G\left(\tfrac{4}{3}\pi R^{3}\rho\right)}{R^{2}}=\tfrac{4}{3}\pi G\rho R .
\]

\stepn{2} \textbf{Read the dependence.} At fixed density, $g\propto R$. Doubling the radius therefore
\textbf{doubles} $g$.
\ans{$g$ doubles}

\smallskip
\textbf{Why this trips people.} The instinct is that $g\propto1/R^{2}$ should make a bigger planet
weaker. But the mass is not held fixed: a bigger planet of the same stuff contains much more of it,
and $R^{3}$ beats $R^{2}$.
\end{exbox}

\begin{exbox}{Escape velocity from the Moon}
The Moon has mass \SI{7.35e22}{kg} and radius \SI{1.74e6}{m}. Find its escape velocity, and comment on
why the Moon has no atmosphere.

\stepn{1} \textbf{Set the total energy to zero}, the condition for just barely escaping:
\[
\tfrac{1}{2}mv_e^{2}-\frac{GMm}{R}=0
\qquad\Longrightarrow\qquad
v_e=\sqrt{\frac{2GM}{R}} .
\]

\stepn{2} \textbf{Substitute.}
\[
v_e=\sqrt{\frac{2(6.67\times10^{-11})(7.35\times10^{22})}{1.74\times10^{6}}}
=\sqrt{5.64\times10^{6}}=\SI{2.37}{km/s}.
\]

\stepn{3} \textbf{Why no atmosphere.} At lunar temperatures a significant fraction of gas molecules
exceed \SI{2.4}{km/s}, so over geological time they leak away, as described in \S\ref{sec:escape}. \ans{\SI{2.37}{km/s}}
\end{exbox}

\begin{exbox}{Raising a satellite's orbit}
A \SI{1000}{kg} satellite is in a circular orbit \SI{400}{km} above the Earth. How much energy is
needed to move it to an orbit of twice that radius? Take $M_E=\SI{6.0e24}{kg}$,
$R_E=\SI{6.4e6}{m}$.

\stepn{1} \textbf{Find the energy of a circular orbit.} Gravity supplies the centripetal force, so
$GMm/r^{2}=mv^{2}/r$, which gives $K=\tfrac12mv^{2}=\dfrac{GMm}{2r}$. Adding $U=-GMm/r$,
\[
E=K+U=-\frac{GMm}{2r}.
\]
(\chref{ch:kepler} draws out its paradox: a higher orbit has more energy but a lower speed.)

\stepn{2} \textbf{Take the difference} between $r_2=2r_1$ and $r_1$:
\[
\Delta E=-\frac{GMm}{2r_2}+\frac{GMm}{2r_1}
=\frac{GMm}{2}\left(\frac{1}{r_1}-\frac{1}{r_2}\right)
=\frac{GMm}{4r_1}.
\]

\stepn{3} \textbf{Substitute} $r_1=6.4\times10^{6}+4\times10^{5}=\SI{6.8e6}{m}$:
\[
\Delta E=\frac{(6.67\times10^{-11})(6.0\times10^{24})(1000)}{4(6.8\times10^{6})}
\approx\SI{1.47e10}{J}=\SI{14.7}{GJ}.
\]
\ans{$\approx\SI{1.5e10}{J}$}
\end{exbox}

\begin{exbox}{Why you weigh less at the equator}
The Earth rotates once in \SI{24}{h}. Estimate how much smaller the apparent weight of a body is at the
equator than at the pole, as a fraction of $mg$. Take $R_E=\SI{6.4e6}{m}$ and $g=\SI{9.8}{m/s^2}$.

\stepn{1} \textbf{Identify the mechanism.} At the pole the body is on the rotation axis and travels no
circle. At the equator it travels a circle of radius $R_E$ once a day, so part of the gravitational
pull is \emph{used up} providing the centripetal force, and the scale reads less.

\stepn{2} \textbf{Write the apparent weight at the equator.} The scale reads the normal force $N$;
gravity $mg$ minus $N$ must supply the centripetal force $m\omega^2R_E$, where $\omega$ is the
Earth's angular speed (\chref{ch:circkin}), so
\[
N=mg-m\omega^{2}R_E .
\]

\stepn{3} \textbf{Evaluate the correction.} With
$\omega=2\pi/(\SI{86400}{s})=\SI{7.27e-5}{rad/s}$,
\[
\omega^{2}R_E=(7.27\times10^{-5})^{2}(6.4\times10^{6})=\SI{0.034}{m/s^2},
\qquad
\frac{\omega^{2}R_E}{g}=\frac{0.034}{9.8}\approx0.0035,
\]
about \textbf{0.35\%}. A \SI{70}{kg} person (weight \SI{686}{N}) weighs about
$70\times0.034\approx\SI{2.4}{N}$ less at the equator.
\ans{$\approx0.35\%$}

\smallskip
\emph{Note:} the measured pole-to-equator difference is larger still, about $0.5\%$, because the Earth
is also flattened, so the equator sits farther from the centre. Rotation causes both effects.
\end{exbox}

\section{Beyond Newton: Einstein's Theory of Gravity}\label{sec:beyondnewton}

\scopebox{This section is not examinable in the Class~11 syllabus and no derivation here is required.
It is included because every equation in this chapter is an \emph{approximation}, and a student who
never learns where the approximation fails will think physics is finished. Read it for orientation;
examine yourself on \S\ref{sec:summary16}.}

\begin{motivbox}{Two things about Newton's law that should bother you}
Newton's law of gravitation is one of the most successful equations ever written. It predicted Neptune
before anyone saw it (\S\ref{sec:newtonlaw}). But two features of it are, on reflection, very strange.

\smallskip
\textbf{First, it is instantaneous.} The formula $F=\dfrac{GMm}{r^{2}}$ contains no time. If the Sun
were to vanish, the equation says the Earth's orbit would straighten \emph{at once}, eight minutes
before anyone saw the Sun go dark. Information would have travelled infinitely fast, which special
relativity forbids: no signal can outrun light.

\smallskip
\textbf{Second, it is an action at a distance.} Two bodies separated by empty space pull on each other
with nothing in between to carry the pull. Newton himself found this intolerable and said so plainly:
he could describe gravity's behaviour but offered no mechanism, writing that he would ``frame no
hypotheses''.

\smallskip
Einstein's general theory of relativity (1915) resolves both. Gravity is not a force transmitted
instantaneously across a gap; it is a \emph{property of spacetime itself}, and changes in it propagate
at exactly the speed of light.
\end{motivbox}

\subsection{The equivalence principle: the idea the whole theory grows from}

General relativity began with a thought experiment that Einstein later called the happiest thought of his
life. Shut in a windowless box, you release a ball and it falls at \SI{9.8}{m/s^2}. Is the box resting on
the Earth, or far out in space and accelerating upward at $g$ (Figure~\ref{fig:equiv})? No experiment inside
the box can tell, and the \emph{equivalence principle} raises this to a law: the effects of gravity and of
acceleration cannot be told apart.

\begin{figure}[htbp]
\centering
\tikzfig{gr_equivalence}
\figcaption{fig:equiv}{Einstein's thought experiment. A sealed box on the Earth's surface and the same
box accelerating at $g$ in deep space give \emph{identical} results for every mechanical experiment. If
no measurement can distinguish them, Einstein argued, then gravity and acceleration are not merely
similar but the same thing.}
\end{figure}

\begin{remarkbox}{Why this is more than a curiosity}
Newtonian mechanics uses the single word ``mass'' for two quite different quantities, and this book has
so far, like most, quietly treated them as one.
\begin{itemize}
\item \textbf{Inertial mass} $m_i$, appearing in $F=m_ia$, measures resistance to being accelerated.
\item \textbf{Gravitational mass} $m_g$, appearing in $F=\dfrac{GMm_g}{r^{2}}$, measures how strongly a
      body responds to gravity.
\end{itemize}
There is no reason in Newtonian physics why these should be equal. Yet experiment finds
$m_i=m_g$ to about one part in $10^{15}$ (the MICROSCOPE satellite, 2022). In Newton's framework this is an unexplained
coincidence; it is the reason all bodies fall at the same rate, a fact so familiar we stop noticing how
odd it is.

\smallskip
Einstein took the coincidence as a \emph{clue}. If gravity and acceleration are genuinely
indistinguishable, then gravity is not a force at all; it is what you feel when you are prevented from
following the natural path through spacetime. A person in free fall feels nothing, not because the
forces cancel, but because there is no force to feel.
\end{remarkbox}

\why{\textbf{The weightless astronaut, again.} \chref{ch:newton} said she floats because she is in free
fall. Einstein's account is that she follows the straightest available path through curved spacetime,
and what you register as ``weight'' standing on the floor is the floor \emph{preventing} you from
following yours.}

\subsection{Gravity as curvature}

If gravity cannot be distinguished from acceleration, then perhaps it is not a force at all. Einstein's answer was
to describe it as the geometry of spacetime (Figure~\ref{fig:curv}).

\begin{figure}[htbp]
\centering
\tikzfig{gr_curvature}
\figcaption{fig:curv}{The standard picture, with the standard warning. A mass distorts the geometry
around it, and other bodies follow the straightest paths available in that distorted geometry. Note that
this diagram is only an analogy: it shows curved \emph{space} on a flat page, whereas the real theory
curves \emph{spacetime}, and it is the time direction that produces most of what we call gravity.}
\end{figure}

\begin{remarkbox}{Einstein's field equation, for orientation only}
The whole theory is contained in one equation relating geometry to matter:
\[
G_{\mu\nu}=\frac{8\pi G}{c^{4}}\,T_{\mu\nu}
\]
The left-hand side describes the curvature of spacetime; the right-hand side describes the energy and
momentum present, and $c$ is the speed of light. You are not expected to use this, or even to read the symbols. The point is the
structure: \textbf{geometry on the left, matter on the right}, with Newton's constant $G$ still sitting
in the bridge between them.

\smallskip
John Wheeler's summary is the one worth remembering:
\begin{center}
\emph{Spacetime tells matter how to move; matter tells spacetime how to curve.}
\end{center}
\end{remarkbox}

\begin{remarkbox}{Newton as the slow, weak limit}
Einstein's theory does not discard Newton's; it \emph{contains} it. When
\begin{itemize}
\item speeds are small compared with $c=\SI{3.00e8}{m/s}$, and
\item gravitational fields are weak, meaning $\dfrac{GM}{rc^{2}}\ll1$,
\end{itemize}
the field equation reduces to $F=\dfrac{GMm}{r^{2}}$, with corrections of relative size about
$\dfrac{GM}{rc^{2}}$. For the Earth's surface $\dfrac{GM}{Rc^{2}}\approx7\times10^{-10}$, so Newton's law
is good to about one part in a billion here, which is why you may use it for every problem in this book.

\smallskip
This is the normal pattern in physics: a deeper theory reproduces the older one where the older one was
tested.
\end{remarkbox}

\subsection{Four places where the difference shows}

Newton's law and Einstein's theory agree to extraordinary precision in everyday conditions. The differences
appear only in strong fields, at high speeds, or in very precise measurements, and four of them are described
below and sketched in Figure~\ref{fig:grev}.

\begin{figure}[htbp]
\centering
\tikzfig{gr_evidence}
\figcaption{fig:grev}{Four effects Newtonian gravity cannot account for. (a) Mercury's perihelion creeps forward, in the same sense as its orbital motion; (b) starlight grazing the Sun is bent towards it, so the star appears displaced outwards; (c) a clock lower in a gravitational well runs slower than a higher one; (d) GPS satellite clocks gain a net \SI{38}{\micro s} a day and must be corrected. The first three are classical tests of general relativity, the fourth an everyday technology that would fail without it. Effects (a)--(c) are greatly exaggerated, and (d) is not to scale.}
\end{figure}

\begin{remarkbox}{The perihelion of Mercury}
By Newton's law, a single planet alone with the Sun would trace an ellipse that closes on itself exactly. Mercury's does not: its
perihelion creeps forward. Most of that creep is caused by the pull of other planets and was
understood, but a residue of
\[
43\ \text{arcseconds per century}
\]
stubbornly resisted explanation for more than half a century; Le~Verrier's Vulcan
(\S\ref{sec:newtonlaw}) was never found, because the missing piece was the law. General relativity predicts the residue exactly, with no free parameters. Einstein
computed it in November 1915, and wrote that the result gave him palpitations.
\end{remarkbox}

\begin{remarkbox}{The bending of starlight}
Light has no mass, so in strict Newtonian gravity it should not be deflected at all. If one nonetheless
treats a light pulse as a fast particle, one obtains a deflection at the Sun's limb of about
$0.87''$ (arcseconds). General relativity, treating light as following the straightest path in curved spacetime,
predicts
\[
\delta=\frac{4GM}{c^{2}R}=1.75''
\]
\emph{exactly twice} the naive value. The expeditions organised by Eddington and Dyson for the solar
eclipse of 1919 found the larger figure, within rather large uncertainties, and Einstein became famous
overnight; radio measurements since have confirmed the factor of two to better than one part in a
thousand. The factor of two is the point: it
distinguishes ``gravity acts on light'' from ``gravity is geometry''.
\end{remarkbox}

\begin{remarkbox}{Gravitational time dilation}
Clocks deeper in a gravitational well run \emph{slower}. To first order, two clocks a height $h$ apart
near the Earth's surface differ in rate by the fraction
\[
\frac{\Delta t}{t}\approx\frac{gh}{c^{2}} .
\]
The effect is tiny (about $2.5\times10^{-15}$ for $h=\SI{22.5}{m}$) but was measured directly in
1959--60 by Pound and Rebka, using a tower just \SI{22.5}{m} high at Harvard. Today's optical atomic
clocks can detect the difference for a height change of a few tens of centimetres: a clock at your feet
runs measurably slower than one at your head.

\smallskip
Note that this follows from the equivalence principle alone, without the full theory: if gravity is
equivalent to acceleration, and an accelerating observer sees clocks at the front run fast, then a clock
higher in a gravitational field must run fast too.
\end{remarkbox}

\begin{remarkbox}{GPS, which would fail without relativity}
This is the one to quote when someone calls relativity impractical. A GPS satellite orbits at about
\SI{20200}{km}, and its clock is affected two ways:
\begin{center}
\begin{tabular}{@{}lrl@{}}
\toprule
Effect & Rate & Cause\\
\midrule
Gravitational (general relativity) & $+\SI{45.7}{\micro s/day}$ & weaker field, clock runs \emph{fast}\\
Motional (special relativity)      & $-\SI{7.2}{\micro s/day}$  & orbital speed, clock runs \emph{slow}\\
\midrule
\textbf{Net}                       & $\mathbf{+\SI{38}{\micro s/day}}$ & \\
\bottomrule
\end{tabular}
\end{center}
Thirty-eight microseconds sounds negligible. But position is inferred from light travel time, and light
covers about \SI{300}{m} in a microsecond, so an uncorrected receiver would accumulate an error of
roughly \textbf{\SI{11}{km} per day}. Every GPS satellite carries the correction in its design.
\end{remarkbox}

\begin{remarkbox}{Where Newton fails outright}
The effects above are small corrections. In three situations the Newtonian picture breaks down
entirely.
\begin{itemize}
\item \textbf{Black holes.} Compress a mass inside its Schwarzschild radius
      $r_s=\dfrac{2GM}{c^{2}}$ and not even light escapes. For the Sun $r_s\approx\SI{2.95}{km}$; for
      the Earth, about \SI{8.9}{mm}. Curiously, the Newtonian escape-velocity argument
      ($v_{\text{esc}}=c$) gives the same radius, but for entirely the wrong reasons.
\item \textbf{Gravitational waves.} Accelerating masses radiate ripples in spacetime. They were predicted
      in 1916 and detected directly a century later, by LIGO in 2015, from two merging black holes.
\item \textbf{Cosmology.} A consistent account of the expanding universe, and of light crossing it,
      requires general relativity.
\end{itemize}
\end{remarkbox}

\begin{trapbox}{What to say, and not say, in an examination}
The Class~11 syllabus is entirely Newtonian, and every problem you will be set should be solved with
$F=\dfrac{GMm}{r^{2}}$ and the results of this chapter. Do not introduce relativity into a numerical
answer. This section tells you \emph{why} the Newtonian results are trustworthy: the condition
$GM/rc^{2}\ll1$.
\end{trapbox}

\section{Summary}\label{sec:summary16}
Every pair of particles attracts with a force $Gm_1m_2/r^2$ along the line joining them. The shell theorem
lets a spherical body be treated as a point at its centre for any point outside it, and makes the field
vanish inside a hollow shell. The potential energy is $-GMm/r$, negative because the pair is bound, and the
sign of the total energy decides whether an orbit is bound, just escapes, or escapes with energy to spare.
In practice the first decision in any problem is whether the point of interest lies outside the body, on its
surface, or inside it: the three regimes obey different laws, and nearly every error in the chapter comes
from applying the wrong one.

\medskip
{\small\renewcommand{\arraystretch}{1.45}
\begin{center}
\begin{tabular}{@{}ll@{}}
\toprule
Quantity & Result\\
\midrule
Universal law        & $F=\dfrac{GMm}{r^{2}}$, attractive, along the line of centres\\[4pt]
Surface gravity      & $g=\dfrac{GM}{R^{2}}=\tfrac{4}{3}\pi G\rho R$\\[4pt]
At height $h$        & $g_h=\dfrac{gR^{2}}{(R+h)^{2}}\approx g\!\left(1-\dfrac{2h}{R}\right)$ for $h\ll R$\\[4pt]
At depth $d$         & $g_d=g\!\left(1-\dfrac{d}{R}\right)$, exactly linear, zero at the centre\\[4pt]
Field                & $\vc{E}_g=\dfrac{\vF}{m}$, units \si{N/kg}, same number as $g$\\[4pt]
Potential energy     & $U=-\dfrac{GMm}{r}$, zero at infinity, negative when bound\\[4pt]
Potential            & $V=-\dfrac{GM}{r}$, a scalar, so it simply adds\\[4pt]
Work to height $h$   & $\Delta U$ from $-GMm/r$; $mgh$ is its $h\ll R$ limit (\S\ref{sec:gpe})\\[4pt]
Escape velocity      & $v_e=\sqrt{\dfrac{2GM}{R}}=\sqrt{2gR}$, independent of the escaping mass\\[4pt]
Orbital (close)      & $v_o=\sqrt{gR}$, so $v_e=\sqrt2\,v_o$\\
\bottomrule
\end{tabular}
\end{center}
}

\medskip
\textbf{Three questions that settle most problems.}
\emph{(1)} Outside, on, or inside? \emph{(2)} Is the quantity a vector (force, field) or a scalar
(energy, potential)? Scalars add arithmetically and are far less work.
\emph{(3)} Is anything conserved? If only gravity acts, total energy is, and for a central force so is
angular momentum.

\medskip
\textbf{Numbers worth carrying:} $g=\SI{9.8}{m/s^2}$; $R_E\approx\SI{6400}{km}$;
$v_e\approx\SI{11.2}{km/s}$; $v_o\approx\SI{7.9}{km/s}$; $G=\SI{6.67e-11}{N.m^2/kg^2}$.

\section{Problem Types}\label{sec:types16}
Five problem types cover the gravitation questions examiners set. The first turns on the contrast between
gravity above and below the surface, where $g$ falls as $1/r^2$ outside but linearly with radius inside;
two turn on escape velocity, its link to orbital speed and what it does and does not depend on; one asks
for the work done against gravity when $mgh$ no longer applies; and one asks how weight and a pendulum's
period change on another planet.

\archq{Variation of $g$ with height or depth}
{``How does $g$ change at height $h$ above, or depth $d$ below, the Earth's surface?''}
{\textbf{Above:} $g_h=\dfrac{gR^2}{(R+h)^2}$, and for $h\ll R$ this gives $g_h\approx g\left(1-\dfrac{2h}{R}\right)$.
\textbf{Below:} only the mass \emph{inside} pulls, giving the exactly linear
$g_d=g\left(1-\dfrac{d}{R}\right)$.}
{Using the ``$1-2h/R$'' approximation when $h$ is comparable to $R$; it is valid only for small
heights. At $h=R$ the exact formula gives $g/4$, whereas the approximation would nonsensically give
$-g$.}
{At a height equal to the Earth's radius $R$ above the surface, the acceleration due to gravity is:}
{$g/4$}{$g/2$}{$g/8$}{zero\archans{A}{$g_h=\dfrac{gR^2}{(R+R)^2}=\dfrac{gR^2}{4R^2}=\dfrac{g}{4}$.
Option (D) comes from wrongly applying the linear depth law $g(1-x/R)$ to a height, and option (B)
from forgetting to square $R+h$.}}

\archq{Escape velocity and orbital speed}
{``Find the escape velocity, or compare it with the orbital speed of a satellite.''}
{Escape: set total energy to zero, $\tfrac{1}{2}mv_e^2=\dfrac{GMm}{R}$, giving
$v_e=\sqrt{\dfrac{2GM}{R}}=\sqrt{2gR}\approx\SI{11.2}{km/s}$. Orbital (close orbit):
$v_o=\sqrt{gR}\approx\SI{7.9}{km/s}$. Hence $\boxed{v_e=\sqrt2\,v_o}$.}
{Applying $v_e=\sqrt2\,v_o$ to speeds at different radii. It compares escape and circular speeds at the
\emph{same} distance from the centre; a satellite high above the surface moves slower than $\sqrt{gR}$.}
{The escape velocity from a planet is $v_e$. The orbital speed of a satellite in a close orbit around it
is:}
{$\sqrt2\,v_e$}{$v_e$}{$v_e/\sqrt2$}{$v_e/2$\archans{C}{$v_e=\sqrt{2gR}$ and $v_o=\sqrt{gR}$, so
$v_o=v_e/\sqrt2\approx0.707\,v_e$.}}

\archq{Work to raise a body far from the surface}
{``How much work is needed to lift a body to height $h$, or to put it into orbit?''}
{Use the difference of $-GMm/r$ between the two radii, not $mgh$ (\S\ref{sec:gpe} shows why $mgh$ is
only its small-height limit). To place the body in a circular orbit there, add its kinetic energy
$\dfrac{GMm}{2(R+h)}$.}
{Writing $mgh$ for a height comparable to $R$. At $h=R$ it overestimates the work by a factor of two.}
{The work needed to raise a body of mass $m$ from the Earth's surface to a height equal to the Earth's
radius $R$ is:}
{$mgR$}{$\tfrac12mgR$}{$\tfrac14mgR$}{$2mgR$\archans{B}{$W=GMm\left(\dfrac1R-\dfrac1{2R}\right)
=\dfrac{GMm}{2R}=\tfrac12mgR$. Option (A) is the $mgh$ answer, which assumes $g$ stays at its surface
value all the way up.}}

\archq{Weight and pendulum on another planet}
{``What is the weight/period on the Moon, where $g$ is one sixth?''}
{Mass never changes; \emph{weight} $=mg$ scales directly with $g$. A pendulum's period
$T=2\pi\sqrt{\ell/g}$ (length $\ell$) scales as $1/\sqrt g$, so weaker gravity means a \emph{longer} period. A
spring--mass oscillator, by contrast, is unaffected: its period does not involve $g$ at all.}
{Saying the mass changes, or that the spring oscillator also slows. Only the pendulum depends on $g$.}
{A pendulum with period $T$ on Earth is taken to the Moon, where $g$ is one sixth. Its new period is:}
{$T/6$}{$T/\sqrt6$}{$\sqrt6\,T$}{$6T$\archans{C}{$T\propto1/\sqrt g$, and $g$ is divided by $6$, so $T$ is
multiplied by $\sqrt6\approx2.45$.}}

\archq{Escape velocity, what it does and does not depend on}
{``Find the escape velocity, or how it changes with the planet's mass, radius or density.''}
{$v_e=\sqrt{\dfrac{2GM}{R}}=\sqrt{2gR}$. Writing $M=\tfrac{4}{3}\pi R^3\rho$ gives
$v_e=R\sqrt{\dfrac{8\pi G\rho}{3}}$, so $v_e\propto R$ at fixed density. It is
\textbf{independent of the escaping body's mass and its direction of projection}.}
{Thinking a heavier projectile needs a bigger escape speed. The projectile's mass cancels on both sides
of the energy equation.}
{The escape velocity from the Earth's surface is about ($g=\SI{9.8}{m/s^2}$, $R=\SI{6400}{km}$):}
{\SI{7.9}{km/s}}{\SI{11.2}{km/s}}{\SI{22.4}{km/s}}{\SI{2.4}{km/s}\archans{B}{
$v_e=\sqrt{2gR}=\sqrt{2(9.8)(6.4\times10^6)}\approx\SI{11.2}{km/s}$. Option (A) is the \emph{orbital}
speed, smaller by $\sqrt2$.}}

\section{Practice Problems}\label{sec:practice16}

The problems below are of standard examination difficulty and are meant to be worked without looking back at the chapter. Answers, with brief working, follow the list.

\begin{enumerate}[itemsep=3pt]
\item At what height above the Earth's surface is $g$ one quarter of its surface value?
\item At what depth is $g$ one half of its surface value?
\item A planet has half the Earth's radius and twice its density. Find $g$ on it as a multiple of
      $g_E$.
\item Find the escape velocity from a planet whose surface gravity is $2g_E$ and whose radius is
      $2R_E$, as a multiple of the Earth's \SI{11.2}{km/s}.
\item How much work is needed to move a \SI{100}{kg} body from the Earth's surface to infinity? Take
      $g=\SI{10}{m/s^2}$ and $R_E=\SI{6.4e6}{m}$.
\item The gravitational potential at a point is \SI{-6.4e7}{J/kg}. Find the escape velocity from that
      point.
\end{enumerate}

\subsection*{Answers}
{\small
\begin{enumerate}[itemsep=2pt]
\item $g_h=g/4$ needs $(R+h)^{2}=4R^{2}$, so $h=R=\SI{6400}{km}$.
\item $g_d=g(1-d/R)=g/2$ gives $d=R/2=\SI{3200}{km}$. Compare the height at which $g$ halves,
$h=0.414R\approx\SI{2650}{km}$ (\S\ref{sec:shell}): near the surface $g$ falls faster going up than
going down.
\item $g=\tfrac{4}{3}\pi G\rho R\propto\rho R$, so $g=(2)(\tfrac{1}{2})g_E=\mathbf{g_E}$: unchanged.
\item $v_e=\sqrt{2gR}$, so $v_e\propto\sqrt{gR}=\sqrt{(2)(2)}=2$ times, that is,
$\mathbf{\SI{22.4}{km/s}}$.
\item $W=\dfrac{GMm}{R}=mgR=100(10)(6.4\times10^{6})=\SI{6.4e9}{J}$.
\item $V=-GM/r$ and $v_e=\sqrt{2GM/r}=\sqrt{2|V|}=\sqrt{2(6.4\times10^{7})}
=\SI{11.3}{km/s}$.
\end{enumerate}
}
